% Generated draft: requires mathematical and source review.
\documentclass[11pt]{article}
\usepackage[margin=25mm]{geometry}
\usepackage{fontspec}
\setmainfont{Latin Modern Roman}
\usepackage{amsmath,amssymb,mathtools,unicode-math}
\setmathfont{Latin Modern Math}
\usepackage{enumitem,needspace,xurl,hyperref}
\setlength{\parindent}{0pt}
\setlength{\parskip}{5pt}
\setlength{\emergencystretch}{3em}
\newcommand{\N}{\mathbb{N}}
\newcommand{\Z}{\mathbb{Z}}
\newcommand{\Q}{\mathbb{Q}}
\newcommand{\R}{\mathbb{R}}
\newcommand{\C}{\mathbb{C}}
\newcommand{\F}{\mathbb{F}}
\newcommand{\A}{\mathbb{A}}
\newcommand{\PP}{\mathbb{P}}
\newcommand{\E}{\mathbb{E}}
\newcommand{\eps}{\varepsilon}
\newcommand{\dd}{\,\mathrm{d}}
\begin{document}
\section*{UnsolvedMath ID: 30004665}
\textbf{Source problem code:} catalogue entry
\subsection{Definitions and mathematical statement}
The supplied record asks for an equivalence involving a string cone \(S_{\mathbf i}\), its defining inequalities, and a function \(W\) restricted to a torus \(T_{\mathbf i}^{\vee}\). It does not specify the underlying algebraic data, the admissible indices \(\mathbf i\), the function \(W\), the torus coordinates, or the particular defining inequalities. The following makes the intended criterion explicit conditionally on those data being supplied.

Fix an admissible choice of the intended data and its designated presentation
\[
S_{\mathbf i}=\{v\in V:\ell_{\alpha}(v)\geq 0\text{ for every }\alpha\in I\},
\]
where \(V\) is the real ambient vector space and the \(\ell_{\alpha}\) are the linear forms in that presentation. The presentation is \emph{nonredundant} if deleting any one inequality strictly enlarges the defined set; equivalently,
\[
\forall\alpha\in I\ \exists v\in V:\quad
\ell_{\alpha}(v)<0
\quad\text{and}\quad
\ell_{\beta}(v)\geq0\ \text{for every }\beta\in I\setminus\{\alpha\}.
\]
Nonredundancy here is a property of the designated presentation, whose identity must be recovered from the source.

Under the intended Laurent-polynomial interpretation, let \(x_1,\ldots,x_n\) be the designated coordinates on \(T_{\mathbf i}^{\vee}\), and write the restriction, after collecting equal monomials, as
\[
\left.W\right|_{T_{\mathbf i}^{\vee}}
 =\sum_{a\in A}c_a x_1^{a_1}\cdots x_n^{a_n},
\qquad A\subseteq\mathbb Z^n\text{ finite},\qquad c_a\neq0.
\]
The exponent condition is
\[
\forall a\in A\ \forall j\in\{1,\ldots,n\}:\quad |a_j|\leq1,
\]
that is, \(A\subseteq\{-1,0,1\}^n\).

\textbf{Conjectural target.} For every admissible choice of the intended source data,
\[
\text{the designated inequalities for }S_{\mathbf i}\text{ are nonredundant}
\quad\Longleftrightarrow\quad
A\subseteq\{-1,0,1\}^n.
\]
This is a conditional reconstruction of the supplied question. The admissibility hypotheses, coordinate choice, Laurent-polynomial interpretation, and correspondence between the restriction and the designated inequalities remain unspecified; the equivalence is not asserted for arbitrary independently chosen cones and Laurent polynomials. The supplied dated assessment reports the exponent condition as sufficient for nonredundancy and leaves the converse conjectural; that assessment has not been independently verified.
\subsection{Short English statement}
The question asks whether a particular defining system for a string cone is nonredundant exactly when every variable has exponent \(-1\), \(0\), or \(1\) in every monomial of the restricted function \(W\). The algebraic setting, coordinates, and defining system must still be identified from the references. Supplied curation labels the problem partially solved and reports sufficiency of the exponent condition, with the converse remaining conjectural.
\subsection{Sources}
\begin{itemize}
\item Supplied local UnsolvedMath registry record, permanent ID \(30004665\), source problem code \texttt{catalogue entry}, entitled \emph{Exponent Criterion for Nonredundant String-Cone Inequalities}; collection: Oberwolfach Reports Open Problems. The record gives proposed year 2021 and names no proposer.
\item Supplied report citation: \emph{Algebraic Groups (2022)}, Oberwolfach Reports, pp.~1087--1148, DOI \texttt{10.4171/owr/2021/20}. \url{https://doi.org/10.4171/owr/2021/20}. The literature-source label is \emph{Oberwolfach Report 20/2021}. The supplied year labels differ and have not been reconciled.
\item Supplied literature reference, attributed to Koshevoy--Schumann and labeled ``redundancy in string-cone inequalities'': \url{https://arxiv.org/abs/2109.14439}.
\item Literature assessment dated exactly \texttt{2026-08-22}: supplied dated curation, not independently verified. It labels the problem \emph{partially solved}, reports that the exponent condition is sufficient for nonredundancy, and describes the converse as conjectural, with only partial cases and evidence reported in the cited work. The record requests independent verification of time-sensitive status and cited claims.
\item The supplied registry labels the statement \texttt{corrected\_verified} and reports a manual semantic, typographic, and LaTeX review claiming self-containedness and preservation of the target. This is supplied review metadata; it does not resolve the missing mathematical specifications identified here. No external source was visited for this draft.
\end{itemize}
\end{document}
